\documentclass[tikz,border=10pt]{standalone}
\usepackage{tikz}

\begin{document}
\begin{tikzpicture}[
  scale=1.15,
  every node/.style={font=\small, align=center}
]
  % расширяем bbox
  \useasboundingbox (-1.2,-1.5) rectangle (10.0,4.2);

  \node at (0,3.2) {\large $T_2$};
  \node at (8.5,3.2) {\large $T_1$};

  \node (F1) at (0,0) {\strut координационные\\нормы};
  \node (F2) at (2.2,0) {\strut ролевые\\статусы};
  \node (F3) at (4.4,0) {\strut деонтические\\силы};
  \node (F4) at (6.6,0) {\strut санкционные\\механизмы};
  \node (F5) at (8.8,0) {\strut условия\\стабильности};

  \draw (0,3.1) -- (F1);
  \draw (0,3.1) -- (F2);
  \draw (0,3.1) -- (F3);
  \draw (8.5,3.1) -- (F1);
  \draw (8.5,3.1) -- (F2);
  \draw (8.5,3.1) -- (F3);
  \draw (8.5,3.1) -- (F4);
  \draw (8.5,3.1) -- (F5);

  \draw (-1,3.9) -- (9.8,3.9);
  \draw (-1,-1.2) -- (9.8,-1.2);
\end{tikzpicture}
\end{document}
